% =====================================================================
% Appendix A --- Equivalences, overlaps and scope of the closure matrix
% ---------------------------------------------------------------------
% REWRITTEN 2026-10-18, v2 (proposal, not yet adopted). Order inverted on
% request: the three scope conditions of the previous version's A.3 open the
% appendix, the coincidences follow, and the algebra that produces them
% closes it. Register rule: no ADR numbers, no run ids, no file, registry or
% generator paths, no \texttt{}. Retains the label app:equivalences, which
% Section 3 cites three times.
% v2 changes vs v1: scope conditions moved to the front; the full-reallocation
% coincidence added as its own row (Section 3 states it as eta = 1).
% =====================================================================

\section{Equivalences, overlaps and scope of the closure matrix}\label{app:equivalences}

\noindent
The evaluation matrix of Section~\ref{sec:application} crosses five
labour-market closures with three financing closures, and several of its
cells turn out to describe one and the same economy. This appendix states the
conditions under which that happens, lists the coincidences with the
mechanism behind each, and closes with the algebra that produces them. The
derivations are in Appendix~\ref{app:proofs}; the model changes that break
the coincidences and restore variety across the matrix are in
Section~\ref{sec:beyond-demand-only}.

\subsection*{Scope: three conditions}

\paragraph{A demand-only shock.} The programme enters final demand at
baseline technology, so the exercise touches neither the price nor the
technology margin. This is the slice on which the whole matrix is evaluated,
and it is the case of our application.

\paragraph{A price block that cannot see demand.} The zero-profit conditions
contain no demand variable and are homogeneous of degree one in prices and
the wage, so relative prices are pinned by technology alone
(Lemma~\ref{lem:price}). Under the consumption price index as numeraire the
real wage then sits at its baseline anchor for every demand vector, every
financing closure and every shock size (Corollary~\ref{cor:realwage}). A
closure whose only task is to move the \emph{real} wage margin is therefore
inert: the margin it is meant to clear is already at rest.

\paragraph{A quantity block that is elasticity-free.} At those prices the
clearing system is exactly affine and independent of the substitution
elasticities $(\theta,\epsilon,\sigma)$ (Lemma~\ref{lem:affine}) --- the
production side of the classical nonsubstitution result
(Remark~\ref{rem:samuelson}). Every closure that pegs the real wage therefore
solves one and the same linear system, whatever the technology.

The third condition is also what locates the boundary of the exercise: all
three hold \emph{on a single common wage}, and Section~\ref{sec:two-questions}
names the single-wage branch and the sectoral family as the matrix's two
labour-market families. Where the wage is common, the conditions above apply
as stated; where it is not, they lapse, and that is where the matrix regains
variety.

\subsection*{The coincidences}

\begin{table}[htbp]
\centering
\footnotesize
\setlength{\tabcolsep}{5pt}
\renewcommand{\arraystretch}{1.25}
\begin{tabular}{@{}>{\raggedright\arraybackslash}p{3.3cm}
                  >{\raggedright\arraybackslash}p{5.5cm}
                  >{\raggedright\arraybackslash}p{2.1cm}
                  >{\raggedright\arraybackslash}p{2.8cm}@{}}
\toprule
Coincidence & Mechanism & Status & Left unidentified \\
\midrule
Fixed real wage $\equiv$ Leontief
& the Leontief corner \emph{is} the fixed-wage closure at the
  zero-substitution limit; a pegged real wage already pins prices, so the
  substitution elasticities have nothing left to move
& exact
& $(\theta,\epsilon,\sigma)$ under a pegged wage \\[3pt]
F2 $\equiv$ F3, mobile rows
& the external transfer is the instrument that holds employment at its pin;
  the two financings differ by exactly the programme cost in that transfer,
  which leaves the real allocation untouched
& exact
& the financing label \\[3pt]
Full reallocation $\equiv$ standard CGE
& at $\eta = 1$ labour is allocated to its cost-minimizing use at a single
  wage, which is precisely the allocation rule the standard CGE closure
  imposes
& definitional
& --- two names for one allocation rule \\[3pt]
Standard CGE $\equiv$ single-wage elastic supply
& with the real wage at its anchor, the supply curve and the full-employment
  pin are the same equation at every supply elasticity and under every
  financing rule
& derivable
& the single-wage supply elasticity \\[3pt]
No reallocation $\equiv$ the rigid corner of the sectoral family
& an infinitely inelastic supply curve reproduces the frozen baseline
  allocation
& exact by construction
& --- it is why no separate rigid sectoral cell is reported \\
\bottomrule
\end{tabular}
\caption{Coincidences within the closure matrix, the mechanism behind each,
their logical status, and the parameter each leaves unidentified.}
\label{tab:equiv-class}
\end{table}

Two qualifications belong with the table. The financing column is the
exception in the fixed-wage rows: there the external transfer is pinned at
zero, so financing \emph{does} bite, and the spread between the tax-financed
and the externally financed variants in those rows is the whole difference
the financing axis carries. And the no-reallocation closure is not a
coincidence but the matrix's one genuinely distinct labour-market regime: it
shares the aggregate employment pin with the standard CGE closure while
freezing the allocation that the latter is free to move, and it is the
endpoint in which the sectoral family terminates.

\subsection*{The algebra behind the coincidences}

\paragraph{The price block.} Unit cost depends on prices and the wage alone,
\begin{equation}
c_i(p,w;A) = \Big( A_i^{\,\epsilon-1}\big( \alpha_i\, w^{1-\epsilon}
+ (1-\alpha_i)\, P_i(p)^{1-\epsilon} \big) \Big)^{\frac{1}{1-\epsilon}},
\qquad
P_i(p) = \Big( \textstyle\sum_u \Omega_{iu}\, p_u^{\,1-\theta} \Big)^{\frac{1}{1-\theta}},
\label{eq:appA-cost}
\end{equation}
and zero profit imposes $p_i = c_i(p,w;A)$ sector by sector. No demand
variable appears in that system, so it is homogeneous of degree one in
$(p,w)$: for any technology, every solution has the form $p = w\,\pi(A)$, and
\emph{relative} prices are functions of technology alone. With the
consumption price index $\Pi(p)$ as numeraire the real wage is therefore a
function of technology alone as well; at baseline technology it sits exactly
at its anchor $w/\Pi(p) = 1$. This is the first scope condition read as
algebra.

\paragraph{The quantity block.} Evaluated at those prices, goods-market
clearing is exactly affine in gross output,
\begin{equation}
y = G\,y + b(R), \qquad
G = \Omega^{\!\top}\mathrm{diag}(a)
  + (1-s)\,\mathrm{diag}(1-m)\,\tilde\omega\,\alpha^{\!\top},
\label{eq:appA-affine}
\end{equation}
and none of the substitution elasticities $(\theta,\epsilon,\sigma)$ enters
$G$ or $b(R)$, where $R$ is the net external transfer to households --- the
variable that clears the external account, and the one the mobile closures
solve for. The technique, and with it relative prices, is independent of the
composition of demand, which allocates quantities only. This is the third
scope condition read as algebra, and it is what makes the matrix a partition
of a few scalar regimes rather than fifteen separate models.

\paragraph{What the algebra does not cover.} Both properties rest on a single
common wage: with a per-sector wage vector the price block is no longer
homogeneous sector by sector, prices become demand-sensitive and both
arguments lapse. Neither property survives a rule that writes the
\emph{nominal} wage as a function of endogenous aggregates --- an aggregate
wage curve, nominal indexation or a re-anchored numeraire --- because such a
rule only re-labels the nominal scale and, under either gauge, is either
vacuous or over-determines the system. A rule on the \emph{real} wage is the
elastic-supply condition already present in the model, and quantity-side
levers such as endogenous saving or endogenous exports move the financing
columns rather than the closures.

What the appendix does not establish is that the matrix is rigid in general.
It is rigid \emph{on the demand-only slice}, and every channel that lets
demand back into the price block --- sectoral labour markets,
demand-sensitive markups, capacity constraints, endogenous import prices,
supply shocks --- restores identification and creates the variety that
Section~\ref{sec:beyond-demand-only} takes up.
